% ============================================================
% C: Encoder and contrastive training   (Anika, ~2:00)
% ============================================================
% Scaffold: the facts below come from PR #1 (src/models/). Check them, replace
% every red TODO, and swap the practice numbers for v0 numbers when you have them.
\section{Method: Model}
\presenter{Anika}

\tikzset{
  cbox/.style={draw=UCFBlack, rounded corners=3pt, fill=LightGray, minimum height=1.0cm,
               align=center, font=\footnotesize, inner sep=4pt},
  cplain/.style={cbox, fill=white},
  chot/.style={cbox, fill=UCFGold!60},
  carr/.style={-{Stealth[length=1.8mm]}, thick, UCFBlack},
  cnote/.style={font=\scriptsize, text=SoftGray, align=center},
}

% ------------------------------------------------------------
\begin{frame}{One encoder for loops and beatboxes}
  \slidepoint{The same network embeds a drum loop and a beatbox, so ``same rhythm'' can become ``nearby vectors''.}

  \centering
  \begin{adjustbox}{max width=\textwidth}
  \begin{tikzpicture}[node distance=0.55cm]
    \node[cplain] (audio) {Audio\\{\scriptsize 22.05 kHz mono}};
    \node[cbox, right=of audio] (fe) {Front-end\\{\scriptsize R1: 3-band onsets}\\{\scriptsize R3: 64-bin log-mel}};
    \node[cbox, right=of fe] (conv) {Stride-2 conv\\{\scriptsize $\sim$20 ms tokens}};
    \node[chot, right=of conv] (tf) {4-layer\\Transformer\\{\scriptsize $d=128$}};
    \node[cbox, right=of tf] (pool) {Attention\\pooling};
    \node[cplain, right=of pool] (out) {128-d\\unit vector};
    \node[cplain, above=0.45cm of conv, font=\scriptsize] (kind) {kind token:\\library / beatbox / tap};
    \draw[carr] (audio) -- (fe); \draw[carr] (fe) -- (conv); \draw[carr] (conv) -- (tf);
    \draw[carr] (tf) -- (pool); \draw[carr] (pool) -- (out);
    \draw[carr] (kind.east) -| (tf.north);
    \node[cnote, below=0.15cm of tf] {weights shared by\\queries and loops};
  \end{tikzpicture}
  \end{adjustbox}

  \vspace{0.35cm}
  \raggedright
  \begin{itemize}
    \item About 550k parameters; trains on one GPU \todo{training time on v0 data}.
    \item Front-ends compared in experiment E2: R1 is timbre-light by design, R3 is a generic spectrogram.
          \todo{which front-end(s) the results use}
  \end{itemize}

  \note{
    \todo{C: speaker notes.}
  }
\end{frame}

% ------------------------------------------------------------
\begin{frame}{Contrastive training}
  \slidepoint{Clips of the same pattern are pulled together, whatever kit or voice plays them; other patterns are pushed apart.}

  \begin{columns}[T,onlytextwidth]
    \column{0.48\textwidth}
      \begin{itemize}
        \item \textbf{Supervised contrastive loss} (SupCon): the class is the pattern ID.
        \item \textbf{Batches} of $P$ patterns $\times$ $K$ views: different kits, different voices, different sloppiness.
              \todo{$P$, $K$, batch size}
        \item AdamW, cosine learning-rate decay, gradient clipping (some seeds collapsed without it).
        \item Validation MRR every epoch; the best checkpoint is kept.
      \end{itemize}
    \column{0.48\textwidth}
      \figtodo[4.8cm]{training curve: loss and validation MRR per epoch, on the v0 data}
  \end{columns}

  \note{
    \todo{C: speaker notes.}
  }
\end{frame}

% ------------------------------------------------------------
\begin{frame}{First check: the pipeline works}
  \slidepoint{On practice data, the encoder ranks the right loop far above chance.}

  \centering
  \begin{tabular}{lcccc}
    \toprule
    Front-end & MRR (95\% CI) & Hit@1 & Hit@5 & Hit@10 \\
    \midrule
    chance & 0.052 & & & \\
    R1: 3-band onsets & 0.373 (0.325--0.424) & 0.22 & 0.52 & 0.74 \\
    R3: log-mel & 0.793 (0.751--0.834) & 0.66 & 0.96 & 0.99 \\
    \bottomrule
  \end{tabular}

  \vspace{0.4cm}
  \raggedright\small
  \begin{itemize}
    \item Practice set: 100 test loops, 200 queries. The ``queries'' were the same loops on a second drum kit,
          not beatboxes, so R3 can partly match by sound. It shows the pipeline works; it doesn't settle R1 vs.\ R3.
    \item \todo{replace with, or add, the v0 numbers on fake beatbox queries}
  \end{itemize}

  \note{
    \todo{C: speaker notes.}
  }
\end{frame}
